\section{欧拉回路与一笔画}
\label{sec:04-Discrete-Mathematics/03-Graph-Theory/07}

18 世纪的 Königsberg 被 Pregel 河切成四块：河心两座岛，南北两岸，七座桥把它们连在一起。城里流传一个消遣：能不能从某处出发，每座桥恰好走一遍，最后回到出发点？

试过的人很多。换一条路线，总是不是漏了一座桥，就是某座桥走了两遍。可``没人走通''离``走不通''还差着一整个论证——尝试再多也只覆盖有限条路线。要断言一件事不可能，穷举帮不上忙。

Euler 在 1736 年的第一步是把绝大部分信息扔掉。陆地多大、桥多长、走得多快、桥是石砌还是木搭，全都不影响答案；剩下的只有``哪块陆地和哪块陆地之间有桥、有几座''。每块陆地缩成一个点，每座桥画成一条线，七桥问题就变成了一张四个顶点、七条边的图。

\begin{center}
\begin{tikzpicture}[x=1cm,y=1cm,font=\small]
  \draw[black!70] (-2.2,0) .. controls (-2.9,1.5) and (-1.5,2.7) .. (0,2.2);
  \draw[black!70] (-2.2,0) .. controls (-1.1,0.4) and (-0.6,1.3) .. (0,2.2);
  \draw[black!70] (-2.2,0) .. controls (-2.9,-1.5) and (-1.5,-2.7) .. (0,-2.2);
  \draw[black!70] (-2.2,0) .. controls (-1.1,-0.4) and (-0.6,-1.3) .. (0,-2.2);
  \draw[black!70] (-2.2,0) -- (2.2,0);
  \draw[black!70] (0,2.2) -- (2.2,0);
  \draw[black!70] (0,-2.2) -- (2.2,0);
  \fill (0,2.2) circle (2.6pt); \node[above] at (0,2.35) {北岸\ \(\deg=3\)};
  \fill (0,-2.2) circle (2.6pt); \node[below] at (0,-2.35) {南岸\ \(\deg=3\)};
  \fill (-2.2,0) circle (2.6pt); \node[left] at (-2.4,0) {岛\ \(\deg=5\)};
  \fill (2.2,0) circle (2.6pt); \node[right] at (2.4,0) {东岛\ \(\deg=3\)};
\end{tikzpicture}

\emph{顶点的地理位置和边的曲直都不再有意义，留下的只有连接方式。四块陆地的度数是 \(5,3,3,3\)，全是奇数——这个事实马上会成为整个论证的支点。}
\end{center}

图上有两条平行边（南岸与岛之间并排两座桥），所以讨论的是多重图；计算度数时平行边各算一次，自环算两次。

\subsection{把``走得通''换算成进出计数}

沿桥行走，每经过一块中途的陆地，必然是沿一座桥进来、沿另一座桥出去。桥被成对地消耗掉。三种情形穷尽了所有可能：中途经过的陆地，每来一次消耗两座桥，度数必须是偶数才能配得完；若起点与终点不同，起点多一次``只出不进''、终点多一次``只进不出''，这两处度数为奇；若行程闭合，最初的离开与最后的返回也配成一对，起点同样是偶度。

Königsberg 的四块陆地度数全奇，哪一种情形都落不进去。问题于是有了确定的否定答案，而这个答案不依赖任何一次具体的尝试，只依赖一次奇偶计数。

Euler 做成的事情比解开一个游戏大得多。``能否遍历所有边''是一个关于整条路线的全局性质，``每个顶点度数的奇偶''是一个逐点就能查完的局部性质，而后者完全刻画了前者。全局被局部彻底刻画，在数学里并不常见，每次出现都值得多看两眼。

\subsection{定理：连通加全偶，不多不少}

经过每条边恰好一次的闭合迹叫欧拉回路，不闭合的叫欧拉通路；中文里更形象的说法是一笔画。\hyperref[sec:04-Discrete-Mathematics/03-Graph-Theory/02]{``路径、环与连通性''}给过判据，这里补上它的来历。

\begin{theorem}
无向图中所有非零度顶点属于同一连通分量时，它存在欧拉回路，当且仅当每个顶点的度数为偶数。
\end{theorem}

条件的措辞需要留心。这里说的不是``图连通''：一个孤立点没有任何关联边，不会妨碍你走遍其余的边，要求它也连进来是多余的约束。如果图里本来就没有孤立点，条件就退回到通常说的``连通且度数全偶''。

必要性就是上面那段计数的形式化。有意思的是充分性，因为它的证明本身就是找路线的办法。

设图连通且全偶。从任一非零度顶点 \(v_0\) 出发，沿没走过的边随便走，走过就划掉。关键在于：除 \(v_0\) 之外，任何顶点你既然能走进去，偶度就保证它还剩至少一条边让你走出来。所以随意行进只可能在回到 \(v_0\) 时停住——只有 \(v_0\) 那次出发没有对应的进入去消耗一条边。你于是拿到一条闭合回路 \(C\)。

如果 \(C\) 用光了所有边，事情就完了。否则删掉 \(C\) 上的边：每个顶点被 \(C\) 经过一次就损失两条关联边，剩余图的度数仍然全偶。连通性保证剩余边中必有一条的端点落在 \(C\) 上。从这个公共顶点出发在剩余图里重走一遍，又得一条闭合回路，把它在那个顶点处嵌进 \(C\)——走到那里时先绕完新回路再接着走 \(C\) 的后半段。合并后的回路更长，度数依旧全偶。反复做下去，边总会用完。

条件在哪一步被用到，可以逐条对账。``全偶''用在两处：保证随意行进不会卡在中途，以及保证删掉一条回路后剩余图仍然全偶，归纳得以继续。``连通''只用在一处：保证剩余的边一定挨着已有的回路，不会剩下一块永远拼不上来的孤岛。缺了前者，行进会死在半路；缺了后者，你会得到几条互不相连的回路，各自完美却拼不成一条。

顺带说明为什么``一直挑没走过的边走到底''还不算算法。想象两个三角形共用一个顶点：全部六条边、五个顶点，度数 \((2,2,4,2,2)\)，全偶且连通，欧拉回路一定存在。可你若先绕完其中一个三角形，就回到了公共顶点、无路可走，另一个三角形原封未动。这不是失败，只是提醒你必须保留``在旧回路的中途插入新回路''这个动作。把这个动作用一个栈实现，就是 Hierholzer 算法：无路可走的顶点才弹出并写进答案，栈里留着的都是还可能被新回路撑开的位置。用邻接表存图，每条边只被删除一次，整个过程是 \(O(\abs{V}+\abs{E})\)，与读入这张图的成本同阶。

这里的``容易''有两层：判据本身简洁到可以口算，而且构造能在输入规模的线性时间内完成。两层都成立的问题不多。

\subsection{不回起点，以及边有方向}

不要求闭合时，进出计数只在两端放宽，于是奇度顶点的个数必须是 0 或 2。握手定理早已保证奇度顶点总数为偶数（见\hyperref[sec:04-Discrete-Mathematics/03-Graph-Theory/01]{``图的语言与度数''}），欧拉通路恰好用尽了这个配额里最小的非零选项。

证明用一个补边的小技巧：在两个奇度顶点 \(s,t\) 之间添一条新边（本来相邻就再加一条平行边），全图变成全偶，取欧拉回路，再把新边删掉，剩下的正是一条从 \(s\) 到 \(t\) 的欧拉通路。而且通路的两端别无选择，只能是这两个奇度顶点。

奇度顶点 0 个，可以一笔画并回到起点，起点随意；恰好 2 个，可以一笔画，但起笔与收笔被钉死；超过 2 个（握手定理下至少 4 个），怎么动笔都不行。Königsberg 有四块奇度陆地，落在最后一类——连``不必回家''的走法都不存在。

两个条件缺一不可，各有各的反例。只查奇偶会漏掉两个互不相连的三角形：每点度数都是 2，奇偶全合格，可一支笔跨不过去。只查连通性则漏掉 Königsberg 本身。

边有方向时，奇偶不再是合适的语言。闭合行程里每沿一条有向边进入，必须沿另一条有向边离开，条件变成

\[
\deg^{+}(v)=\deg^{-}(v)
\qquad\text{对每个顶点 }v,
\]

再加上所有含边顶点属于同一个强连通分量。开放的有向欧拉通路允许两端失衡：起点多一次离开，终点多一次进入。无向图里的``两个奇度顶点''被换成``一处多出一次离开、一处多出一次进入''，骨架没变——给每一次进入找到一次配对的离开。

\subsection{判据失败之后}

只要必须逐一覆盖的对象是边，欧拉模型就是正确的第一站：街道巡检、铲雪路线、电路板钻孔走线、管网巡查。检查连通性、数一遍奇度顶点，答案几乎立刻出来，用不着搜索。

现实里更常见的是判据失败。邮递员必须走遍辖区每条街道，而图里偏偏有奇度顶点，重复走路无法避免。问题于是从``能不能不重复''变成``怎样重复才最省''。这就是管梅谷 1962 年提出的中国邮路问题。核心观察很短：把一条边多走一遍，等于在它两端之间添一条平行边；要让全图变成全偶，复制的那些边必须恰好把奇度顶点两两配对。于是最优方案落在``求一组奇度顶点的配对，使配对两端的最短路长度之和最小''，奇偶计数再一次把走法问题化成了图上的纯计算。

模型的边界同样值得看清。边如果带长度或代价，重复路段的成本不是重数乘单位长度那么简单；道路如果有方向，必须换用入度出度平衡的版本；图如果在变（修路、通桥），一次判据的结论不能顺延到变化之后的图。

最要紧的一条边界不在这些工程细节里，而在问题的措辞上。欧拉模型要求不可遗漏的对象是边。如果任务是拜访每一个客户，而客户坐落在顶点上，那么度数奇偶不但给不出答案，连问的问题都不是同一个。下一节就把``每条边一次''换成``每个顶点一次''，看看这一个词的差别能把计算难度推到多远。
